\documentclass[11pt]{article}

\usepackage[T1]{fontenc}
\usepackage{lmodern}
\usepackage[a4paper,margin=29mm]{geometry}
\usepackage{amsmath,amssymb,amsthm,mathtools}
\usepackage{microtype}
\setlength{\emergencystretch}{2em}
\usepackage{enumitem}
\usepackage{xcolor}
\usepackage[colorlinks=true,linkcolor=blue!45!black,citecolor=blue!45!black,
  urlcolor=blue!55!black]{hyperref}

\newtheorem{theorem}{Theorem}[section]
\newtheorem{proposition}[theorem]{Proposition}
\newtheorem{lemma}[theorem]{Lemma}
\newtheorem{corollary}[theorem]{Corollary}
\newtheorem{definition}[theorem]{Definition}
\newtheorem{remark}[theorem]{Remark}

\newcommand{\C}{\mathbb C}
\newcommand{\R}{\mathbb R}
\newcommand{\N}{\mathbb N}
\newcommand{\XiR}{\Xi}
\newcommand{\cZ}{\mathcal Z}
\newcommand{\cM}{\mathcal M}
\newcommand{\cQ}{\mathcal Q}
\newcommand{\cR}{\mathcal R}
\newcommand{\ord}{\operatorname{ord}}
\newcommand{\Res}{\operatorname{Res}}
\newcommand{\supp}{\operatorname{supp}}

\title{Fixed-Window Logarithmic Weil Tomography}
\author{Michael Haufschild \and GPT 5.6 Sol \and Fable 5 \and Grok 4.5}
\date{}

\begin{document}

\maketitle

\begin{abstract}
For every fixed width $L>0$ we introduce one translated rectangular two-edge family in the
centered Weil form of the Riemann xi-function.  Its window transform
\[
 A_L(w)=\frac{1-e^{-Lw}}{w}
\]
has no zeros in either open transverse half-plane.  We prove that logarithmic sampling of the
translation parameter, $a=\frac12\log n$, turns the growing edge into an ordinary Dirichlet
carrier.  The resulting single scalar sequence is a complete transverse detector: the Riemann
hypothesis is equivalent to its nonnegativity, to its eventual lower boundedness, and to its norm
boundedness.  If the hypothesis is false, the sequence is cofinally unbounded below, and its sharp
power exponent is the supremal horizontal displacement of the xi divisor.

Our principal theorem is a cutoff-free multiplicity-exact tomography statement.  The Dirichlet
series of the complete sequence of actual Weil values admits a canonical meromorphic continuation
$\cM_L$ to $\Re s>1$.  Its simple poles are exactly
\[
  s=1+(\rho-\tfrac12),\qquad \XiR(\rho)=0,\quad \Re\rho>\tfrac12.
\]
If $\rho$ has analytic multiplicity $m(\rho)$, then
\[
  \operatorname*{Res}_{s=\rho+1/2}\cM_L(s)
  =2m(\rho)A_L(\rho-\tfrac12)^2.
\]
Since the window factor never vanishes, one fixed boundary window and one arithmetic mesh recover
the complete off-critical divisor with multiplicity, not merely its support or rightmost displacement.
We also identify the sequence with a fixed-width second
difference of the xi screw function.  The results are unconditional divisor-side theorems; the
remaining RH-bearing problem is to construct this completed observable from the
prime--archimedean or semilocal operator realization and prove a cofinal lower floor or holomorphic
right-half-plane continuation there.
\end{abstract}

\noindent\textbf{2020 Mathematics Subject Classification.}
11M26, 11M06, 30B50, 42A82, 47B25.

\noindent\textbf{Keywords.}
Riemann hypothesis, Weil quadratic form, screw function, Dirichlet series, Landau theorem,
Connes trace formula, pole tomography.

\section{Introduction}

Weil's quadratic-form criterion converts the Riemann hypothesis into positivity of a Hermitian
form on a test space~\cite{Weil1952,Yoshida1992,Bombieri2000}.  Connes's trace-formula program
places the same form in a spectral and semilocal operator architecture~\cite{Connes1999}, with
later work isolating the scaling Hamiltonian, the archimedean positivity mechanism, prolate
near-radical directions, and zeta cycles~\cite{ConnesConsani2019,ConnesConsani2021,
ConnesConsani2023}.  Recent finite spectral triples package prime cutoffs into self-adjoint
finite objects, while leaving convergence to the completed xi divisor as the decisive global
step~\cite{ConnesConsaniMoscovici2025}.

The usual Weil criterion ranges over a large test space.  This is natural for positivity but
obscures a more economical question: can a single explicit cofinal family already carry every
transverse zero?  One-parameter criteria are known in screw-function and Hilbert-space form.
Suzuki proved several screw-function equivalents of RH~\cite{Suzuki2023} and later obtained a
screw-line equality criterion on a cofinal interval~\cite{Suzuki2025}.  His recent synthesis
places the Yoshida, Bombieri, Connes, and Connes--Consani--Moscovici forms in one screw-function
framework~\cite{Suzuki2026}.  Finite Galerkin values have also been identified exactly with
finite Guinand--Weil zero sums~\cite{Groskin2026}.

There is also a classical superzeta literature in which Dirichlet series are formed directly over
the Riemann zeros and their meromorphic polar structure is studied~\cite{Voros2003}.  That
construction takes the zero set as the Dirichlet indexing object.  Here the Dirichlet variable is
the ordinary integer $n$: the coefficients are one logarithmically sampled Weil sequence, and its
Mellin poles recover the transverse divisor.  This distinction is essential for the intended
prime-side and semilocal interpretation.

This paper isolates a different compression.  Fix any $L>0$ once and for all.  Translate the
same two-edge rectangle to aperture $a$, then sample only at $a=\frac12\log n$.  Three phenomena
occur simultaneously.

\begin{enumerate}[label=(\roman*)]
\item The rectangle has an exact half-plane nonvanishing law, so no generic profile selection is
needed.
\item The cofinal translation becomes an ordinary Dirichlet monomial $n^w$ at a centered zero
$w=\rho-\frac12$.
\item The discarded edge and the on-line divisor remain uniformly bounded, so the transverse
Dirichlet carrier controls the entire Weil sequence.
\end{enumerate}

The principal theorem attaches a cutoff-free meromorphic continuation to the complete actual
scalar sequence and proves
multiplicity-exact divisor tomography: pole positions give the shifted right xi support and
normalized residues give the analytic multiplicities.  Functional-equation and Schwarz symmetry
then recover the full off-critical divisor.  Companion theorems give three exact RH criteria,
cofinal divergence under false RH, the sharp transverse exponent, and the exact screw-coordinate
identity.

The novelty is the conjunction of fixed positive width, arithmetic logarithmic sampling,
cofinal divergence, sharp transverse exponent, and observable multiplicity-exact pole tomography.  Landau's
theorem, Weil positivity, shifted-zeta continuation, and screw functions are the parent
mechanisms.  We keep this boundary explicit.  We also keep the logical boundary explicit: the
paper does not prove RH.  It identifies the exact cofinal scalar carrier on which a successful
prime-side or semilocal construction must establish a completed source--divisor coupling and a
cofinal lower floor or holomorphic continuation.

\section{The centered xi divisor and the rectangular window}

Let
\[
  \XiR(s)=\frac12s(s-1)\pi^{-s/2}\Gamma(s/2)\zeta(s).
\]
Its nontrivial zeros form a multiset invariant under $\rho\mapsto1-\rho$ and
$\rho\mapsto\overline\rho$.  Write
\[
  w=\rho-\frac12,
\]
and let $\cZ$ denote the resulting centered multiset, with analytic multiplicity.  Then
\[
  |\Re w|<\frac12,qquad
  \sum_{w\in\cZ}\frac1{|w|^2}<\infty.
  \tag{2.1}
\]
The latter follows from the classical zero count $N_\Xi(T)=O(T\log T)$.

For $L>0$ define
\[
  A_L(w)=
  \begin{cases}
  \dfrac{1-e^{-Lw}}{w},&w\ne0,\\[6pt]
  L,&w=0,
  \end{cases}
  \qquad
  F_{a,L}(w)=e^{aw}A_L(w)+e^{-aw}A_L(-w).
  \tag{2.2}
\]
The transform is the bilateral centered Laplace transform of
\[
  \mathbf 1_{[a-L,a]}+\mathbf 1_{[-a,-a+L]}.
\]
We work directly with the absolutely convergent divisor polarization
\[
  Q_L(a)=\sum_{w\in\cZ}F_{a,L}(w)^2.
  \tag{2.3}
\]
Absolute convergence follows from (2.1) and the strip bound below.  Formula (2.3) is the
rectangle-completed zero side of the centered Weil form.  Its unconditional Hermitian symmetry
makes $Q_L(a)$ real.  No distributional assertion about a
discontinuous test function is needed for the arguments in this paper.

\begin{lemma}[Half-plane nonvanishing]\label{lem:nonvanishing}
If $L>0$ and $\Re w\ne0$, then $A_L(w)\ne0$.
\end{lemma}

\begin{proof}
If $A_L(w)=0$ with $w\ne0$, then $e^{-Lw}=1$.  Taking absolute values gives
$e^{-L\Re w}=1$, hence $\Re w=0$, a contradiction.
\end{proof}

\begin{lemma}[Exact edge decomposition]\label{lem:edge}
For every $w\in\C$,
\[
  A_L(-w)=e^{Lw}A_L(w),
  \tag{2.4}
\]
and consequently
\[
  F_{a,L}(w)^2=A_L(w)^2
  \bigl(e^{2aw}+2e^{Lw}+e^{-2(a-L)w}\bigr).
  \tag{2.5}
\]
Moreover, uniformly on $|\Re w|\le\frac12$,
\[
  |A_L(w)|\le \frac{1+e^{L/2}}{|w|}\qquad(w\ne0).
  \tag{2.6}
\]
\end{lemma}

\begin{proof}
Both identities are algebraic consequences of (2.2).  For (2.6), use
$|1-e^{-Lw}|\le1+e^{-L\Re w}\le1+e^{L/2}$.
\end{proof}

Let $\cZ_+=\{w\in\cZ:\Re w>0\}$, retaining multiplicity, and define
\[
  c_L(w)=A_L(w)^2,
  \qquad
  d_L(n)=2\sum_{w\in\cZ_+}c_L(w)n^w
  \quad(n\ge1).
  \tag{2.7}
\]
Conjugation invariance makes $d_L(n)$ real.  Functional-equation symmetry folds the growing edge
of (2.5) onto (2.7).

\begin{proposition}[Uniform remainder]\label{prop:remainder}
For each fixed $L>0$ there is a finite constant $K_L$ such that, whenever
$a\ge L$,
\[
  \left|Q_L(a)-D_L(a)\right|\le K_L,
  \tag{2.8}
\]
where $D_L(\frac12\log n)=d_L(n)$.  One may take
\[
  K_L=R_L\sum_{w\in\cZ}|w|^{-2}
  \tag{2.9}
\]
with an explicit $R_L$ depending only on $L$.
\end{proposition}

\begin{proof}
For $\Re w>0$, select the first term in (2.5); for $\Re w<0$, select the third
term.  The functional equation identifies the latter with the first term at $-w$.  The middle
term is independent of $a$.  On the selected positive half-strip the third term has modulus at
most a constant depending only on $L$ when $a\ge L$; the symmetric statement holds on the
negative half-strip.  On the critical line all exponential factors in (2.5) have bounded modulus.
Lemma~\ref{lem:edge} therefore bounds every discarded summand by $R_L|w|^{-2}$.
Summation using (2.1) gives (2.8).
\end{proof}

\section{The logarithmic Dirichlet carrier}

Assume $\cZ_+\ne\varnothing$ and set
\[
  \delta_* = \sup_{w\in\cZ_+}\Re w.
  \tag{3.1}
\]
Then $0<\delta_*\le\frac12$.  Define
\[
  C_L(s)=\sum_{w\in\cZ_+}c_L(w)\zeta(s-w)
  \tag{3.2}
\]
initially on $\Re s>1+\delta_*$.  Absolute convergence permits interchange of the two sums and
gives
\[
  \sum_{n\ge1}\frac{d_L(n)}{n^s}=2C_L(s).
  \tag{3.3}
\]

The continuation of (3.2) requires a multiplicity-safe formulation.  Put
\[
  X_L(s)=\XiR(s-\tfrac12)C_L(s)
  \tag{3.4}
\]
on the initial half-plane, filling each removable shifted-zeta singularity.  The unconditional
zero count yields
\[
  \sum_{w\in\cZ_+}|w|^{-3/2}<\infty.
  \tag{3.5}
\]
Together with a square-root vertical bound for $\zeta$ on closed substrips, (3.5) gives local
normal convergence of the filled series on $\Re s>1$.  Thus $X_L$ is holomorphic there and
\[
  \widetilde C_L(s)=\frac{X_L(s)}{\XiR(s-\tfrac12)}
  \tag{3.6}
\]
is a canonical meromorphic continuation of $C_L$.

\begin{proposition}[Exact pole support]\label{prop:poles}
For $L>0$, the meromorphic function $\widetilde C_L$ has, in $\Re s>1$, a simple pole at
$s=1+w$ for every $w\in\cZ_+$, and no other poles.  If $w$ has multiplicity $m$, then
\[
  \begin{aligned}
  \ord_{1+w}X_L&=m-1,
  &\ord_{1+w}\XiR(s-\tfrac12)&=m,\\
  \ord_{1+w}\widetilde C_L&=-1,
  &\operatorname*{Res}_{s=1+w}\widetilde C_L(s)&=m c_L(w).
  \end{aligned}
  \tag{3.7}
\]
\end{proposition}

\begin{proof}
Fix a distinct support point $w_0\in\cZ_+$ of multiplicity $m$.  Near $p=1+w_0$, factor
\[
  \XiR(s-\tfrac12)=(s-p)^m u(s),\qquad u(p)\ne0.
  \tag{3.8}
\]
The $m$ equal-support summands in (3.4) are identical.  The pole-filled factor
$\XiR(s-\frac12)\zeta(s-w_0)$ has exact order $m-1$ at $p$; multiplying by
$c_L(w_0)\ne0$ and by the positive integer $m$ does not change that order.  Every summand with
support different from $w_0$ is divisible by $(s-p)^m$.  Local normal convergence preserves this
common order floor through their infinite sum.  Hence the equal-support fiber has order $m-1$,
the off-support tail has order at least $m$, and their sum has order exactly $m-1$.  Division by
(3.8) proves the order statement.  More precisely, put
$R(z)=(z-1)\zeta(z)$ with its removable value $R(1)=1$.  The equal-support fiber is
\[
  m c_L(w_0)(s-p)^{m-1}u(s)R(s-w_0),
\]
whereas the off-support tail is divisible by $(s-p)^m$.  Dividing by (3.8) gives principal part
$m c_L(w_0)/(s-p)$ and proves the residue formula in (3.7).  Away from the coordinates $1+w$,
the denominator in (3.6) is nonzero and the
quotient is holomorphic.
\end{proof}

\begin{proposition}[Pole bound from tail summability]\label{prop:polebound}
Fix $L>0$, $C\in\R$, $N\in\N$, and $\delta\ge0$ such that
$\Re w\le\delta$ for every $w\in\cZ_+$.  Put
\[
 b_n=\begin{cases}d_L(n)+C,&n\ge N,\\0,&n<N.\end{cases}
\]
If the ordinary Dirichlet series of $(b_n)$ converges at a real point $\sigma>1$, then
\[
  \Re(1+w)\le\sigma\qquad(w\in\cZ_+).
  \tag{3.9}
\]
\end{proposition}

\begin{proof}
Convergence at the real point $\sigma$ gives holomorphy of the ordinary L-series on
$\Re s>\sigma$.  On a farther right half-plane it agrees with the canonical continuation built
from $\widetilde C_L$, the constant-zeta term, and a finite Dirichlet polynomial.  The identity
theorem propagates the agreement through the connected pole-free part of $\Re s>\sigma$.
If some shifted pole $p=1+w$ lay in that half-plane, the ordinary L-series would give a
holomorphic continuation across $p$.  Multiplication by $\XiR(s-\frac12)$ would then force the
cancelled numerator to vanish to order at least the multiplicity $m$ of $w$.  Proposition
\ref{prop:poles} gives exact order $m-1$, a contradiction.  Hence no such $p$ exists.
\end{proof}

\begin{remark}
The order computation is the point at which multiplicity matters.  Treating only simple zeros
would hide the finite equal-support fiber and would not justify the absence of cancellation.
\end{remark}

\section{Cofinal divergence and the RH criteria}

We use the classical Landau principle: an ordinary Dirichlet series with eventually nonnegative
real coefficients cannot be holomorphic through the real point at its abscissa of convergence.
The form needed below is scale-local: if the series converges at one real point and extends
holomorphically to a disc reaching a smaller real point, its Taylor expansion forces convergence
at that smaller point.  This is the standard proof of Landau's theorem applied after deleting a
finite prefix.

\begin{theorem}[Cofinal fixed-window completeness]\label{thm:completeness}
Fix $L>0$ and put
\[
  q_L(n)=Q_L\!\left(\frac12\log n\right)
  \qquad\left(L\le\frac12\log n\right).
  \tag{4.1}
\]
Then the following are equivalent.
\begin{enumerate}[label=(\alph*)]
\item The Riemann hypothesis holds.
\item $q_L(n)\ge0$ at every admissible logarithmic sample.
\item The admissible tail of $q_L(n)$ is bounded below.
\item The admissible values $|q_L(n)|$ are bounded.
\end{enumerate}
If RH is false, then for every $B\in\R$ and every $N_0$ there is an admissible $n\ge N_0$ such
that
\[
  q_L(n)<B.
  \tag{4.2}
\]
\end{theorem}

\begin{proof}
Under RH every $w\in\cZ$ is purely imaginary.  The two-edge transform $F_{a,L}(w)$ is then real
in the spectral coordinate, so every summand in (2.3) is a real square.  This proves (a)$\Rightarrow$(b).
The uniform reciprocal-square estimate also gives (a)$\Rightarrow$(d), while (b)$\Rightarrow$(c)
and (d)$\Rightarrow$(c) are immediate.

Suppose RH is false.  Functional-equation symmetry gives $\cZ_+\ne\varnothing$.  If $d_L(n)$
were eventually bounded below, add a constant and delete a finite prefix to obtain a nonnegative
coefficient sequence.  Its L-series agrees on $\Re s>1+\delta_*$ with
\[
  2C_L(s)+C\zeta(s)-P(s),
  \tag{4.3}
\]
where $P$ is a finite Dirichlet polynomial.  Let $\sigma_c$ be its abscissa of convergence.
Initially $\sigma_c\le1+\delta_*$.  The shifted-zeta carrier is holomorphic on a disc centred at
the real boundary point $1+\delta_*$: the right divisor has a uniform positive ordinate gap and
the reciprocal-$3/2$ majorant gives local normal convergence there.  Landau's theorem therefore
excludes $\sigma_c=1+\delta_*$.  Thus the nonnegative tail series converges at some real
$\sigma<1+\delta_*$.  Proposition~\ref{prop:polebound} gives
$1+\delta_*\le\sigma$, a contradiction.  Therefore $d_L(n)$ is cofinally below every real
level.  Proposition~\ref{prop:remainder} transfers the same statement to $q_L(n)$, proving
(4.2).  Consequently (c) implies RH, and all four conditions are equivalent.
\end{proof}

The same argument retains quantitative information.

\begin{theorem}[Sharp transverse exponent]\label{thm:exponent}
Assume RH is false and let $\delta_*$ be given by (3.1).  There is a constant $M_L<\infty$ such
that every admissible sample satisfies
\[
  |q_L(n)|\le M_L n^{\delta_*}.
  \tag{4.4}
\]
Conversely, for every $0\le\alpha<\delta_*$, every $C\ge0$, and every $N_0$, there is an
admissible $n\ge N_0$ for which
\[
  |q_L(n)|>Cn^\alpha.
  \tag{4.5}
\]
Thus $\delta_*$ is the sharp power exponent of the single fixed-window sequence.
\end{theorem}

\begin{proof}
By (2.6), $\sum_{w\in\cZ_+}|c_L(w)|<\infty$.  Since
$|n^w|=n^{\Re w}\le n^{\delta_*}$, (2.7) gives the upper bound for $d_L(n)$; the bounded
remainder gives (4.4).

If (4.5) failed for some $\alpha<\delta_*$, then the bounded remainder would give an eventual
$O(n^\alpha)$ bound for $d_L(n)$.  Its L-series would converge absolutely at every
$\Re s>1+\alpha$.  Choose $\beta$ with $\alpha<\beta<\delta_*$.  The continuation argument of
Proposition~\ref{prop:poles}, now applied at the real line $1+\beta$, would force every shifted
pole to satisfy $\Re(1+w)\le1+\beta$.  Taking suprema yields $\delta_*\le\beta$, a contradiction.
The bounded difference between $d_L$ and $q_L$ is absorbed into $Cn^\alpha$ because
$n^\alpha\ge1$.
\end{proof}

\section{Cutoff-free multiplicity-exact divisor tomography}

Proposition~\ref{prop:poles} concerns the isolated growing carrier.  We now remove every cutoff
from the final object.  Define the global observable remainder
\[
  e_L(n)=q_L(n)-d_L(n).
  \tag{5.1}
\]
Choose an auxiliary $N$ so that $n\ge N$ implies $n\ge1$ and
$L\le\frac12\log n$.  Proposition~\ref{prop:remainder} gives
\[
  |e_L(n)|\le K_L\qquad(n\ge N).
  \tag{5.2}
\]
The terms with $n<N$ form a finite sequence.  Consequently
\[
  E_L(s)=\sum_{n\ge1}\frac{e_L(n)}{n^s}
  \tag{5.3}
\]
converges absolutely and is holomorphic on $\Re s>1$.  Although a cutoff proves this fact, neither
$e_L$, $E_L$, nor the continuation below depends on it.

\begin{theorem}[Main theorem: cutoff-free observable divisor tomography]\label{thm:tomography}
For every $L>0$, the ordinary Dirichlet series of the complete observable sequence
\[
  \sum_{n\ge1}\frac{q_L(n)}{n^s}
  \tag{5.4}
\]
has, on its initial half-plane of convergence, the canonical continuation
\[
  \cM_L(s)=2\widetilde C_L(s)+E_L(s).
  \tag{5.5}
\]
The function $\cM_L$ is meromorphic on $\Re s>1$, and
\[
  \ord_s\cM_L=-1
  \quad\Longleftrightarrow\quad
  s=1+w\text{ for some }w\in\cZ_+.
  \tag{5.6}
\]
At every other point of $\Re s>1$, $\cM_L$ is holomorphic.  Moreover, if
$w\in\cZ_+$ has analytic multiplicity $m(w)$, then
\[
  \operatorname*{Res}_{s=1+w}\cM_L(s)
  =2m(w)c_L(w)=2m(w)A_L(w)^2,
  \tag{5.7}
\]
and therefore
\[
  m(w)=
  \frac{\operatorname*{Res}_{s=1+w}\cM_L(s)}{2A_L(w)^2}.
  \tag{5.8}
\]
Thus one full observable sequence, with no deleted prefix and no auxiliary choice, determines the
right xi divisor with its exact analytic multiplicities.
\end{theorem}

\begin{proof}
The coefficient identity $q_L=d_L+e_L$, absolute convergence, and (3.3) give (5.5) on
$\Re s>1+\delta_*$.  The correction $E_L$ is holomorphic on $\Re s>1$.  At every shifted right
zero, Proposition~\ref{prop:poles} gives order $-1$ for $2\widetilde C_L$; addition of a
holomorphic function cannot cancel a pole of strictly smaller order.  Away from those shifted
coordinates both terms in (5.5) are holomorphic.  This proves (5.6).  The correction has zero
residue at $1+w$, while Proposition~\ref{prop:poles} gives residue $m(w)c_L(w)$ for
$\widetilde C_L$.  The factor two gives (5.7).  Lemma~\ref{lem:edge} shows $A_L(w)\ne0$, so
division yields (5.8).
\end{proof}

\begin{corollary}[Divisor recovery and RH]\label{cor:tomography}
The complete right off-critical divisor of xi, including multiplicities, is recovered from one
complete fixed-width sequence.
Its support is
\[
  \{\rho:\XiR(\rho)=0,\ \Re\rho>\tfrac12\}
  =\{s-\tfrac12:\Re s>1,\ \ord_s\cM_L=-1\}.
  \tag{5.9}
\]
The multiplicity at $\rho=s-\tfrac12$ is the normalized residue in (5.8).  Functional-equation
and Schwarz symmetry then recover the entire off-critical divisor.  In particular, RH is
equivalent to holomorphy of $\cM_L$ on $\Re s>1$.
\end{corollary}

\begin{corollary}[Intrinsic normalized pole divisor]\label{cor:intrinsic}
Let $L_1,L_2>0$ and let $w$ be a right off-critical support point.  Then
\[
 \frac{\operatorname*{Res}_{s=1+w}\cM_{L_1}(s)}{2A_{L_1}(w)^2}
 =
 \frac{\operatorname*{Res}_{s=1+w}\cM_{L_2}(s)}{2A_{L_2}(w)^2}
 =m(w).
 \tag{5.10}
\]
Thus the normalized weighted pole divisor is independent of the rectangular width.  There is no
cutoff coordinate to quotient out.
\end{corollary}

\begin{proof}
Apply (5.7) to each $L_j$.
\end{proof}

\begin{remark}
The continuation is not an arbitrary interpolation of the scalar samples.  It is fixed by the
xi-cancelled shifted-zeta carrier and the globally defined absolutely convergent remainder
L-series.  This is the incidence marking that prevents a
transverse relocation of the divisor while leaving only ordinate data unchanged.
\end{remark}

\section{The screw-function coordinate}

Define the centered xi screw profile
\[
  \Psi_\Xi(t)=\sum_{w\in\cZ}\frac{\cosh(tw)-1}{w^2},
  \qquad t\in\R.
  \tag{6.1}
\]
The series converges absolutely by (2.1), because $|\Re w|<\frac12$, and it is real by divisor
symmetry.

\begin{theorem}[Exact fixed-width curvature identity]\label{thm:screw}
For every $L,a\in\R$,
\[
  Q_L(a)=2\Bigl[
    \Psi_\Xi(2a)+\Psi_\Xi(2a-2L)-2\Psi_\Xi(2a-L)+2\Psi_\Xi(L)
  \Bigr].
  \tag{6.2}
\]
Consequently Theorem~\ref{thm:completeness} is equivalently an RH criterion for the one
fixed-width second-difference expression in (6.2), sampled only at $2a=\log n$.
\end{theorem}

\begin{proof}
For $w\ne0$, substitute $A_L(w)=(1-e^{-Lw})/w$ and
$A_L(-w)=e^{Lw}A_L(w)$ into $F_{a,L}(w)^2$.  Expanding the exponentials and grouping each
$e^{tw}+e^{-tw}-2=2(\cosh(tw)-1)$ gives the summand identity corresponding to (6.2).
Absolute convergence permits summation over $\cZ$.  The removable value at $w=0$ follows by
continuity, although $w=0$ is not a xi zero.
\end{proof}

The screw identity places the theorem directly beside the modern screw-line formulation of the
Weil form~\cite{Suzuki2023,Suzuki2025,Suzuki2026}.  The additional feature here is the exact
rectangular stencil and the arithmetic logarithmic mesh, which produce the Dirichlet continuation
and pole tomography of Section~5.

\section{Implications for the Connes route}

The theorem changes the positive obligation on a semilocal or prime-side construction.  It is not
necessary to establish positivity on an unspecified compact family, nor to match finitely many
spectral ordinates.  For any chosen $L>0$, it is enough to construct the actual completed
source--divisor incidence along the single section $q_L(n)$ and prove one cofinal finite lower floor
\[
  q_L(n)\ge-C\qquad(n\ge N).
  \tag{7.1}
\]
This closes RH by Theorem~\ref{thm:completeness}; norm boundedness and positivity are stronger
sufficient special cases.  Alternatively, the same source construction may prove directly that
the cutoff-free canonical continuation $\cM_L$ is holomorphic on $\Re s>1$.

This is a cofinal, noncompact carrier condition.  Finite-prime positivity or a finite
self-adjoint spectrum does not imply (7.1): the construction must transport the completed
Euler--Gamma incidence into the complete actual sequence.  The compatible finite-factorization
viewpoint is therefore useful only when its cofinal section controls the global remainder in
(5.3).  Once
that happens, Theorem~\ref{thm:tomography} gives a rigid alternative:

\begin{quote}
the residual Mellin channel is holomorphic on $\Re s>1$, or its weighted simple-pole divisor is
exactly the shifted off-critical xi divisor, with multiplicities encoded by normalized residues.
\end{quote}

In particular, the noncompact escape required by prolate or near-radical Connes--Consani
mechanisms is not an undirected loss of compactness.  Along the rectangular section it must be a
coherent escape whose scalar trace values obey (7.1), or whose Mellin channel is holomorphic on the
whole target half-plane.  The same obligation applies to recent
finite-prime spectral triples.  Their normalized determinants must carry the completed divisor
strongly enough to identify the full continuation, not only a finite list of real spectral zeros.

\section{Scope and limitations}

The results are unconditional, but they are divisor-side theorems.  They do not derive the lower
bound (7.1) from primes, Gamma factors, a Sonin space, or a self-adjoint operator.  They therefore
do not prove RH.  Their role is to isolate a minimal positive carrier and to show exactly what it
remembers.

The rectangle is deliberately nonsmooth.  We avoid an inadmissible appeal to the classical smooth
test space by defining $Q_L$ through the absolutely convergent xi divisor polarization.  A
prime-side explicit-formula realization of this exact boundary value requires either an
appropriate completion theorem or a justified approximation whose limit preserves the scalar
sequence.  That source theorem is the remaining RH-bearing step.

Every off-critical zero, simple or multiple, produces a simple pole, so pole order alone forgets
multiplicity.  The main theorem repairs exactly that loss: the pole residue is the analytic
multiplicity times the explicit nonvanishing rectangular weight.  Thus the weighted pole divisor,
not the unweighted pole set, is the intrinsic observable.

Finally, the novelty claim is intentionally narrow.  Weil positivity, the Connes trace-formula
architecture, Landau's theorem, and the screw function are established parent theories.  The new
result is the fixed-window logarithmic compression and its quantitative and meromorphic
consequences.

\begin{thebibliography}{99}

\bibitem{Bombieri2000}
E.~Bombieri,
\emph{Remarks on Weil's quadratic functional in the theory of prime numbers, I},
Rend. Lincei Mat. Appl. \textbf{11} (2000), 183--233.

\bibitem{Connes1999}
A.~Connes,
\emph{Trace formula in noncommutative geometry and the zeros of the Riemann zeta function},
Selecta Math. (N.S.) \textbf{5} (1999), 29--106.

\bibitem{ConnesConsani2019}
A.~Connes and C.~Consani,
\emph{The scaling Hamiltonian},
arXiv:1910.14368, 2019.

\bibitem{ConnesConsani2021}
A.~Connes and C.~Consani,
\emph{Weil positivity and trace formula, the archimedean place},
arXiv:2006.13771, revised 2026.

\bibitem{ConnesConsani2023}
A.~Connes and C.~Consani,
\emph{Spectral triples and zeta-cycles},
Enseign. Math. \textbf{69} (2023), 93--148.

\bibitem{ConnesConsaniMoscovici2025}
A.~Connes, C.~Consani, and H.~Moscovici,
\emph{Zeta spectral triples},
arXiv:2511.22755, 2025.

\bibitem{Groskin2026}
A.~Groskin,
\emph{A finite Guinand--Weil dictionary and archimedean tail order for the truncated Weil
quadratic form},
arXiv:2607.02828, 2026.

\bibitem{Maurizi2010}
B.~Maurizi,
\emph{Extending Landau's theorem on Dirichlet series with non-negative coefficients},
arXiv:1009.0228, 2010.

\bibitem{Suzuki2023}
M.~Suzuki,
\emph{Aspects of the screw function corresponding to the Riemann zeta function},
arXiv:2206.03682v4, 2023.

\bibitem{Suzuki2025}
M.~Suzuki,
\emph{On the Hilbert space derived from the Weil distribution},
Canad. J. Math., published online 3 November 2025,
doi:10.4153/S0008414X25101739.

\bibitem{Suzuki2026}
M.~Suzuki,
\emph{Weil's quadratic form via the screw function},
arXiv:2606.09096, 2026.

\bibitem{Titchmarsh1986}
E.~C.~Titchmarsh,
\emph{The Theory of the Riemann Zeta-Function}, 2nd ed., revised by D.~R.~Heath-Brown,
Oxford University Press, 1986.

\bibitem{Voros2003}
A.~Voros,
\emph{Zeta functions for the Riemann zeros},
Ann. Inst. Fourier (Grenoble) \textbf{53} (2003), 665--699;
erratum \textbf{54} (2004), 1139.

\bibitem{Weil1952}
A.~Weil,
\emph{Sur les ``formules explicites'' de la th\'eorie des nombres premiers},
Commun. S\'em. Math. Univ. Lund, volume dedicated to Marcel Riesz (1952), 252--265.

\bibitem{Yoshida1992}
H.~Yoshida,
\emph{On Hermitian forms attached to zeta functions},
in \emph{Zeta Functions in Geometry}, Adv. Stud. Pure Math. \textbf{21},
Mathematical Society of Japan, 1992, 281--325.

\end{thebibliography}

\end{document}
